\subsection{\texorpdfstring{CONTINUOUS HOMOMORPHISMS OF \(R^{m}\) INTO \(T^{n}\)}{CONTINUOUS HOMOMORPHISMS OF R TO THE M INTO T TO THE N}}

PROPOSITION 3. Every continuous homomorphism of \(\mathbf{R}^m\) into \(\mathbf{T}^n\) is of the form \(x \to \varphi(u(x))\) , where \(\varphi\) is the canonical homomorphism of \(\mathbf{R}^n\) onto \(\mathbf{T}^n\) (identified with \(\mathbf{R}^n / \mathbf{Z}^n\) ) and \(u\) is a linear mapping of \(\mathbf{R}^m\) into \(\mathbf{R}^n\) .

Let \(f\) be a continuous homomorphism of \(\mathbf{R}^m\) into \(\mathbf{T}^n\) . We shall show that there is a linear mapping \(u\) of \(\mathbf{R}^m\) into \(\mathbf{R}^n\) such that the homomorphisms \(x \to f(x)\) and \(x \to \varphi(u(x))\) coincide at all points of a neighbourhood of \(o\) in \(\mathbf{R}^m\) ; the proposition will then follow from Corollary 1 to Proposition 2 of no. 2. Now let \(V\) be a neighbourhood of \(o\) in \(\mathbf{R}^n\) such that \(\varphi\) , restricted to \(V\) , is a local isomorphism of \(\mathbf{R}^n\) with \(T^n\) , and let \(\psi\) be the inverse of \(\varphi\) , defined on \(\varphi(V)\) . Since \(f\) is continuous, \(V' = \overline{f}^{-1}(\varphi(V))\) is a neighbourhood of \(o\) in \(\mathbf{R}^m\) ; the mapping

\[
\boldsymbol {x} \rightarrow \psi (f (\boldsymbol {x})),
\]

restricted to \(V'\) , is a continuous mapping of \(V'\) into \(R^n\) such that \(\psi(f(x+y)) = \psi(f(x)) + \psi(f(y))\) for each pair of points x, y in \(R^m\) such that \(x \in V'\) , \(y \in V'\) and \(x + y \in V'\) ; hence (no. 2, Corollary 1 to Proposition 2) this mapping coincides with a well-determined continuous homomorphism u of \(R^m\) into \(R^n\) at all points of a neighbourhood W of o in \(\mathbf{R}^m\) . By Proposition 1 (no. 1) \(u\) is a linear mapping of \(\mathbf{R}^m\) into \(\mathbf{R}^n\) ; for all \(x \in W\) we have therefore \(f(x) = \varphi(u(x))\) , which completes the proof.

Remark. The same argument shows, more generally, that if \(\varphi\) is a strict morphism of \(R^{n}\) into a group G, whose restriction to a suitable neighbourhood of o is a local isomorphism of \(R^{n}\) with G, then every continuous homomorphism of \(R^{n}\) into G is of the form \(x \to \varphi(u(x))\) , where u is a linear mapping of \(R^{m}\) into \(R^{n}\) .

In the case \(m = n = 1\) , Proposition 3 gives:

PROPOSITION 4. If \(\varphi\) is the canonical homomorphism of \(\mathbf{R}\) onto \(\mathbf{T}\) , every continuous homomorphism of \(\mathbf{R}\) into \(\mathbf{T}\) is of the form \(x \to \varphi(ax)\) where \(a \in \mathbf{R}\) ; and it is a strict morphism of \(\mathbf{R}\) onto \(\mathbf{T}\) if \(a \neq 0\) .

